%%%%%%%%%%%%
%
% $Autor: Wings $
% $Datum: 2019-03-05 08:03:15Z $
% $Pfad: TemplateSensor $
% $Version: 4250 $
% !TeX spellcheck = en_GB/de_DE
% !TeX encoding = utf8
% !TeX root = filename 
% !TeX TXS-program:bibliography = txs:///biber
%
%%%%%%%%%%%%

% Structure
\chapter{Drehpotentiometer}

%Notizen:

%Drehpotentiometer
%Widerstand 10 kOhm, linear
%Achse 6 mm 
%Spannungsfestigkeit: 250 VDC	
%Befestigung Gewinde M10
%
%{../../Images/BOM/Drehpotentiometer_10kOhm_linear_6mm.jpg} 
%
%https://www.reichelt.de/de/de/shop/produkt/drehpotentiometer_10_kohm_linear_6_mm-232554

\section{General}

General description

cite books

\section{Specific Sensor/Actor}

cite board

\section{Specification}

\begin{itemize}
  \item Cite the data sheet
  \item Extract the information from the data sheet
  \item Circuit Diagram
\end{itemize}

\section{Library}

\subsection{Description}

\subsection{Installation}

\begin{itemize}
    \item data types
    \item data structure
    \item data quantity
    \item data quality
\end{itemize}

\subsection{Functions}

\section{Example}

In this section, a simple example is described in detail, which demonstrates how the library supports the sensor/actor.


\subsection{Manual}

\subsection{Inside the Example}

\subsection{Code}

\subsection{Files}



\section{Calibration}

cite method

\section{Simple Code}


\section{Simple Application}



\section{Tests}

\subsection{Simple Function Test}

\subsection{Test all Functions}



\section{Further Readings}

